\documentclass[12pt]{article}

\usepackage[spanish]{babel}
\usepackage[utf8]{inputenc}
\usepackage[T1]{fontenc}
\usepackage{setspace}
\usepackage{geometry}
\usepackage{titlesec}
\usepackage{indentfirst}
\usepackage{apacite}

\geometry{letterpaper, margin=2.54cm}

% Normas APA
\setstretch{2}
\setlength{\parindent}{1.27cm}

% Quitar numeración de secciones
\titleformat{\section}
{\bfseries\large\centering}
{}
{0em}
{}

\begin{document}

\begin{center}
\textbf{\LARGE ÁRBOL DE PROBLEMAS} \\[0.5cm]
\textbf{METODOLOGÍA DE INVESTIGACIÓN} \\[0.5cm]
\textbf{INGENIERÍA DE SISTEMAS} \\[1cm]

ARIANA MARCELA ANDRADE BELLO \\
EMMANUEL ESTEBAN RESTREPO PATARROYO \\[1cm]

POPAYÁN \\
MAYO 2026
\end{center}

\vspace{1cm}

\section{Planteamiento del problema}

En la actualidad, las plataformas digitales y las redes sociales se han convertido en uno de los principales medios de acceso a la información para gran parte de la población. Sin embargo, junto con este crecimiento también ha aumentado la circulación de noticias falsas o información no verificada \cite{shu2017fake}.

Este fenómeno puede generar desinformación, influir en la opinión pública y afectar la forma en que las personas interpretan diferentes acontecimientos sociales, políticos o económicos. Además, la rápida propagación de este tipo de contenido dificulta su control y aumenta su impacto en la sociedad \cite{shu2017fake}. 

Frente a este panorama, resulta relevante investigar soluciones tecnológicas capaces de identificar y analizar la veracidad de la información que circula en entornos digitales antes de que esta sea difundida masivamente \cite{zhou2020survey}. 

Un sistema de este tipo debería ser capaz de examinar el contenido de las noticias, contrastar la información con diferentes fuentes confiables y detectar posibles inconsistencias o contradicciones \cite{castillo2011information}. 

Aunque se han desarrollado diferentes soluciones basadas en inteligencia artificial para detectar noticias falsas, todavía existen varias limitaciones importantes. La mayoría de estos sistemas analizan únicamente el contenido del texto y utilizan datos previamente entrenados, lo que hace que funcionen bien en ciertos casos, pero tengan dificultades para adaptarse a información nueva o a contextos reales \cite{thorne2018automated}.  

Además, muchos de estos enfoques no comparan la información con otras fuentes, por lo que una noticia puede ser clasificada sin verificar si realmente coincide o contradice lo que otros medios están reportando \cite{popat2018declare}.  

En este sentido, la principal diferencia de esta propuesta es que utiliza modelos basados en arquitecturas Transformer, combinados con técnicas de recuperación de información y comparación semántica, tales como BERT, SBERT y enfoques tipo Retrieval-Augmented Generation (RAG), con el fin de contrastar el contenido de una noticia con múltiples fuentes y detectar inconsistencias \cite{ghadiri2022automated}. 

A partir de este proceso de comparación, el sistema buscaría identificar posibles contradicciones o inconsistencias en los datos presentados. Como resultado, se generaría un indicador de confiabilidad que permita orientar a los usuarios sobre la posible veracidad de la información \cite{dey2024retrieval}. 

\section{Pregunta Problema}

¿Cómo influye un modelo basado en inteligencia artificial que permita analizar el contenido de noticias digitales y contrastarlo con múltiples fuentes de información en la detección de posibles inconsistencias y en la determinación de su nivel de confiabilidad?

\section{Objetivo General}

Proponer un modelo basado en inteligencia artificial que analice el contenido de noticias digitales y lo contraste con múltiples fuentes de información para detectar inconsistencias y determinar su nivel de confiabilidad.

\section{Objetivos Específicos}

\begin{itemize}
    \item Analizar las características del contenido de noticias digitales y los criterios de confiabilidad utilizados en distintas fuentes de información.
    
    \item Desarrollar un modelo basado en inteligencia artificial que permita comparar noticias digitales con múltiples fuentes para identificar posibles inconsistencias.
    
    \item Evaluar el desempeño del modelo propuesto en la detección de inconsistencias y en la determinación del nivel de confiabilidad de las noticias digitales.
\end{itemize}

\section{Cadena de Búsqueda}

\begin{verbatim}
("artificial intelligence" OR AI OR "machine learning" OR NLP)
AND ("fake news" OR misinformation OR disinformation)
AND ("fact checking" OR "credibility assessment")
AND (reliability OR trustworthiness OR inconsistencies)
AND ("multiple sources" OR "cross-source analysis")
\end{verbatim}

\vspace{1cm}

\bibliographystyle{apacite}
\bibliography{bibliography}

\end{document}